%% ==========================================================================
%% DRAFT -- August 4, 2026. Author byline: Daniel Kirtchakov (Independent
%% researcher, Half Ounce Research).
%%
%% Claims were re-scoped on 2026-08-04 after an adversarial prior-art sweep
%% and an independent re-verification of the certificate corpus. Sources read
%% in full from raw form on that date:
%%   * arXiv:2308.07915v2 (Bravyi, Cross, Gambetta, Maslov, Rall, Yoder),
%%     main + supplemental LaTeX. Table 3 lists [[288,12,18]] WITHOUT the
%%     "<=" that marks upper-bound-only entries; the "unlikely to be tight"
%%     caveat in the main text is about the CIRCUIT-level distance d_circ.
%%   * arXiv:2605.16523v1 (LEAN-QEC), all 10 source files.
%%     ALSO: repository VerifiedQC/Lean-QEC at commit c73827d (2026-07-10),
%%     whose notes report a completed bv_decide verification of the
%%     [[144,12,12]] distance. The repository has moved past the paper; we
%%     therefore claim NO priority for a machine-checked gross-code distance.
%%   * arXiv:2606.12445 (Chen, Jafari, Lai), main.tex + result tables.
%%     Their [[288,12,18]] table reports d >= 11 (solver-asserted, 7200 s);
%%     repository guluchen/QDistSAT contains no proof artifacts.
%%   * arXiv:2606.02418 (Cruz-Benito, Cross, Kremer, Faro; IBM, 1 Jun 2026):
%%     the gross-code distance d = 12 confirmed exactly by MILP at MIP gap 0.
%%     No certificate is emitted.
%%   * arXiv:2603.22532 (Webster, Jacob, Higgott), full text.
%% Independent audit: qec/INDEPENDENT-VERIFICATION.md, 47/47 checks passed,
%%   5.08 GiB LRAT replayed in pure Python, all five BB matrices rebuilt
%%   byte-identically from arXiv:2308.07915. Four discrepancies (D1-D4) are
%%   reported verbatim in Section 8 of this note; D2 has since been closed.
%% The dated sweep record is SWEEP-RECORD-QEC-2026-08-04.md, released with
%%   the artifacts.
%% ==========================================================================
\documentclass[11pt]{amsart}

\usepackage[margin=1.15in]{geometry}
\usepackage{amsmath,amssymb}
\usepackage{booktabs}
\usepackage{array}
\usepackage[colorlinks=true,allcolors=blue]{hyperref}

\newcolumntype{L}[1]{>{\raggedright\arraybackslash}p{#1}}
\emergencystretch=3em
\tolerance=1200
\hbadness=4000

\newcommand{\FF}{\mathbb{F}_{2}}
\newcommand{\rowsp}{\operatorname{rowsp}}
\newcommand{\rank}{\operatorname{rank}}
\newcommand{\wt}{\operatorname{wt}}
\newcommand{\supp}{\operatorname{supp}}
\newcommand{\HX}{H_{X}}
\newcommand{\HZ}{H_{Z}}
\newcommand{\dX}{d_{X}}
\newcommand{\dZ}{d_{Z}}
\newcommand{\dcirc}{d_{\mathrm{circ}}}
\newcommand{\gross}{[[144,12,12]]}

\theoremstyle{plain}
\newtheorem{theorem}{Theorem}[section]
\newtheorem{lemma}[theorem]{Lemma}
\newtheorem{proposition}[theorem]{Proposition}
\newtheorem{corollary}[theorem]{Corollary}
\theoremstyle{definition}
\newtheorem{definition}[theorem]{Definition}
\theoremstyle{remark}
\newtheorem{remark}[theorem]{Remark}

\begin{document}

\title[Replayable distance certificates for stabilizer codes]{Replayable
minimum-distance certificates for stabilizer codes, with no solver and no
proof assistant in the trusted base: the bivariate-bicycle family through
$n=288$}

\author{Daniel Kirtchakov}
\address{Independent researcher, Half Ounce Research}
\email{daniel@halfounce.io}

\thanks{\emph{Computation and authorship.} All searches, encodings,
symmetry-breaking constructions, and certificate designs in this work were
produced by \textbf{Claude Fable~5} (Anthropic), directed by the author, on a
single Apple~M4 laptop (10 cores, 16\,GB). The three independent checkers
share no code with the generating pipeline, import only the Python standard
library, and were exercised against the corpus by a separate agent instance
that was given no access to the pipeline. This is a factual methods statement,
and it is part of the point of the paper: the artifacts are designed so that
the provenance of the \emph{search} is irrelevant to the validity of the
\emph{result}. External tools used by the pipeline only: CaDiCaL~3.0.1
(\texttt{--lrat --no-binary}), Python~3.}

\thanks{\emph{Prior-art record.} The credit statements in
\S\ref{ssec:credit} are the output of a full-text adversarial sweep completed
on August~4, 2026, in which every source was read in raw form (arXiv LaTeX
source, GitHub repository trees at a named commit, package documentation) and
in which two claims of an earlier internal draft were found to be wrong and
withdrawn. The dated record of that sweep, listing what was read and what
changed, ships with the artifacts as
\texttt{SWEEP-RECORD-QEC-2026-08-04.md}. \emph{Draft of August 4, 2026.}}

\subjclass[2020]{Primary 81P73; Secondary 94B05, 68V15, 03B35}
\keywords{Quantum error correction, minimum distance, stabilizer codes,
bivariate bicycle codes, gross code, LRAT, proof certificates, SAT}

\date{August 4, 2026}

\begin{abstract}
This is a verification contribution, not a discovery. The minimum distances of
the codes treated here are already in the literature; what did not exist is a
standalone artifact that a third party can replay without a SAT solver and
without a proof assistant. We supply one for eleven stabilizer codes, together
with three short checkers, and we submit the result to an independent audit:
$47$ certificate checks run, $47$ passed, $0$ failed; $5.08$\,GiB of LRAT
replayed in pure Python inside a $79$\,MB peak resident set; all five
bivariate-bicycle parity-check matrices rebuilt \emph{byte-identically} from
the published construction of Bravyi et al.; $182$ of $182$ manifest SHA-256
entries matching; $11$ of $11$ negative controls correctly rejected.

The formats are deliberately small. Upper bounds are witness pairs verified by
three $\FF$ matrix--vector products. Lower bounds are LRAT unsatisfiability
proofs for a canonical CNF encoding of ``some nontrivial logical has weight at
most $K$'', in which the checker regenerates the CNF from the raw parity-check
matrices and machine-checks the three algebraic side conditions that make the
encoding \emph{exact}; neither the solver nor the shipped CNF is trusted, and
for every code in the corpus, including the largest proof, replay needs
nothing beyond the Python standard library.

For the gross code $\gross$ of Bravyi et al.\ we certify $d=12$ end to end:
weight-$12$ $X$- and $Z$-witnesses, and in each sector a single
\emph{symmetry-free} LRAT proof ($868$\,MB for $X$, $672$\,MB for $Z$) of
$\dX,\dZ\ge 12$, so that no symmetry lemma enters the trusted base for this
code. For $[[288,12,18]]$ we certify $\dX\ge 14$ by a $2.94$\,GB LRAT proof,
replayed in pure Python in $414$\,s. Bravyi et al.\ assert $d=18$ for that
code exactly, by integer programming and without a checkable artifact; the
strongest lower bound previously \emph{reported} by a certifying-capable
method is $d\ge11$, solver-asserted by Chen, Jafari and Lai with no proof
files in their repository. Our $\dX\ge14$ improves on that bound and is, as
far as our sweep could determine, the only machine-checkable lower bound on
record for this code. It is not evidence against $d=18$.

We state the trusted base explicitly, separating what is machine-checked per
certificate from the facts that are assumed, and we report the four defects
the audit found, including one latent soundness hole in a checker branch.
We claim priority neither for the distance values, which are Bravyi et al.'s,
nor for machine-checked quantum distance proofs, which are LEAN-QEC's and
whose public repository reports a completed gross-code verification as of
2026-07-10. What is offered here is a certificate format and a trusted base:
a $649$-line standard-library reader, and artifacts that outlive the tools
that produced them.
\end{abstract}

\maketitle

\section{Introduction}\label{sec:intro}

\subsection{The problem}
The minimum distance $d$ of a quantum error-correcting code is the single
number that determines how many errors it can detect. It is also expensive:
computing it is a minimum-weight problem over a coset space, intractable in
general \cite{Vardy97}, and the standard tools for it in the quantum
literature are of two kinds. Randomized searches such as QDistRnd
\cite{QDistRnd} and Stim's \texttt{search\_for\_undetectable\_logical\_errors}
\cite{Stim} return upper bounds and say so plainly; Stim's documentation is
explicit that ``\textsc{this is a heuristic method}''. Exact methods ---
integer programming \cite{LandahlAndersonRice}, SAT and MaxSAT
\cite{ChenJafariLai}, branch and bound --- return a number and an implicit
assertion that the search was exhaustive. The survey of Webster, Jacob and
Higgott \cite{WJH} classifies the landscape carefully; what almost none of it
produces is an artifact a third party can check.

This matters more than usual here, because the published distances of the
codes now being built into hardware are load-bearing engineering constants.
The bivariate-bicycle (BB) codes of Bravyi, Cross, Gambetta, Maslov, Rall and
Yoder \cite{BCGMRY} --- in particular the $\gross$ ``gross code'' --- are the
basis of a fault-tolerant memory architecture. Their distances were computed
by the mixed-integer-programming method of Landahl, Anderson and Rice
\cite{LandahlAndersonRice}; the paper is explicit about this and about which
entries of its Table~3 are upper bounds only. The gross-code value has since
been confirmed independently, again by mathematical programming: Cruz-Benito,
Cross, Kremer and Faro \cite{CruzBenito} close the MILP to a gap of zero and
obtain $d=12$ exactly. Nothing is wrong with either computation. But a reader
who wants to check them must re-run a solver and trust the re-run; neither
emits an object that survives the solver.

\subsection{What this note provides}
An artifact-first alternative. For each code we ship files that a skeptic can
replay in isolation, and three short Python programs that replay them. The
design constraints were:

\begin{enumerate}
\item[(i)] \emph{No solver in the trusted base.} The checker never runs a SAT
solver and never reads the CNF file the solver saw. It \emph{regenerates} the
CNF from the raw parity-check matrices, then replays the LRAT proof against
its own clause list. A corrupted shipped CNF changes nothing --- this was
tested (\S\ref{sec:audit}).
\item[(ii)] \emph{No proof assistant in the trusted base.} No Lean, no
Mathlib, no toolchain fetch. The reader needs CPython and the files.
\item[(iii)] \emph{The encoding's exactness is machine-checked, not asserted.}
The reduction from ``$\dX\ge K+1$'' to ``this CNF is unsatisfiable'' is valid
only under three algebraic side conditions on $(\HX,\HZ,P)$. The checker
verifies all three, per certificate, over the raw matrices
(Theorem~\ref{thm:exact}).
\item[(iv)] \emph{Everything the checker assumes is written down.}
Section~\ref{sec:tcb} is the honest inventory, including the items argued in
prose and not machine-verified anywhere.
\end{enumerate}

\subsection{Results}\label{ssec:results}
Table~\ref{tab:main} is the corpus. Every entry traces to a file in
\texttt{certificates/}; the byte counts are the actual proof sizes and the
replay times are wall clock from the independent audit of
\S\ref{sec:audit}, measured with \texttt{/usr/bin/python3} (CPython 3.9.6, no
\texttt{numpy}, no compiled helper of any kind).

\begin{table}[t]
\small
\begin{tabular}{lllrrr}
\toprule
code & $n,k$ & certified & lower-bd.\ proof & solver & pure-Python \\
     &       &           & (bytes)          & (s)    & replay (s)  \\
\midrule
Steane $[[7,1,3]]$        & $7,1$   & $d=3$ & $699 \mid 699$ & $0.02\mid0.07$ & $0.0003\mid0.0009$\\
five-qubit $[[5,1,3]]$    & $5,1$   & $d=3$ & $1{,}514$ & $0.02$ & $0.0004$\\
rot.\ surface $d{=}3$     & $9,1$   & $d=3$ & $518 \mid 746$ & $0.03\mid0.02$ & $0.0004\mid0.0004$\\
rot.\ surface $d{=}5$     & $25,1$  & $d=5$ & $4{,}685 \mid 7{,}085$ & $0.02\mid0.02$ & $0.001\mid0.001$\\
rot.\ surface $d{=}7$     & $49,1$  & $d=7$ & $24{,}939 \mid 75{,}107$ & $0.02\mid0.02$ & $0.005\mid0.009$\\
Golay $[[23,1,7]]$        & $23,1$  & $d=7$ & $226{,}377$ each & $0.04\mid0.04$ & $0.024\mid0.024$\\
BB $[[72,12,6]]$          & $72,12$ & $d=6$ & $1.54\mid1.93$\,MB & $0.26\mid0.30$ & $0.14\mid0.17$\\
BB $[[90,8,10]]$          & $90,8$  & $d=10$ & $128\mid151$\,MB & $48.4\mid50.8$ & $10.9\mid13.0$\\
BB $[[108,8,10]]$         & $108,8$ & $d=10$ & $71.6\mid82.3$\,MB & $28.2\mid24.7$ & $6.2\mid7.1$\\
\textbf{BB} $\gross$      & $144,12$& $\mathbf{d=12}$ & $868\mid672$\,MB$^{\dagger}$ & $342\mid227$ & $176\mid72.9$\\
BB $[[288,12,18]]$        & $288,12$& $14\le d\le 18$ & $2.94$\,GB$^{\ddagger}$ & $513$ & $414$\\
\bottomrule
\end{tabular}
\medskip
\caption{The certificate corpus. Paired entries are $X\mid Z$ sector. Solver
times are CaDiCaL 3.0.1 on one Apple M4 laptop, recorded in each code's
\texttt{meta.json}. Replay times are from the independent audit
(\S\ref{sec:audit}) and use nothing outside the Python standard library.
$^{\dagger}$Symmetry-\emph{free} single-instance proofs; a symmetry-broken
$X$ certificate ($124$\,MB${}+{}$$34$\,kB, solver $45.0$\,s) is also shipped
and replays in $10.1$\,s. $^{\ddagger}$Symmetry-broken, $K=13$, two instances
($2{,}941{,}958{,}076 + 191{,}479$ bytes); the ladder rungs $\dX\ge10$
($65$\,MB) and $\dX\ge12$ ($351$\,MB) are shipped as well. At $n=288$ the
certified quantity is $\dX$; the passage to $d$ uses the shipped duality
certificate together with $d=\min(\dX,\dZ)$, and that certificate was
generated after the audit closed (Remark~\ref{rem:no288dual}).}
\label{tab:main}
\end{table}

Three things in that table deserve to be stated precisely, because each of
them touches published work.

\smallskip
\noindent\textbf{(A) The gross code.} We certify $d(\gross)=12$ end to end.
The upper bound is a pair of weight-$12$ witnesses, $970$-byte JSON files,
checked in $7$\,ms each. The lower bound is certified twice in each direction:
once by a symmetry-broken pair of instances in the $X$ sector, and once, in
\emph{each} sector, by a single LRAT proof with no symmetry hypothesis at all
--- $867{,}803{,}294$ bytes for $X$ and $671{,}988{,}205$ bytes for $Z$. The
symmetry-free pair is the one that matters for the trusted base: it removes
the orbit lemma of \S\ref{ssec:sym} from the argument entirely, so
$d(\gross)=12$ rests on nothing but the exactness theorem, the LRAT semantics,
and the standard cardinality and XOR encodings.

We do \emph{not} claim to be first to a machine-checked distance for this
code. The value $d=12$ is Bravyi et al.'s \cite[Table 3]{BCGMRY}, confirmed
exactly at MIP gap zero by Cruz-Benito, Cross, Kremer and Faro
\cite{CruzBenito} and reproduced by SAT in \cite{ChenJafariLai}. The idea of
a kernel-checked quantum distance proof is LEAN-QEC's \cite{LeanQEC}, and
their public repository reports the gross code completed; see
\S\ref{ssec:credit} for the dated record. What is ours is the certificate
format, the symmetry-free variant, and the trusted base.

\smallskip
\noindent\textbf{(B) $[[288,12,18]]$.} We certify $\dX\ge 14$ by a
$2.94$\,GB LRAT proof and $\dX\le 18$ by a weight-$18$ witness. The value is
not in dispute and we are not disputing it: Bravyi et al.\ assert $d=18$ for
this code exactly, by integer programming, and their Table~3 marks it without
the ``$\le$'' that flags their upper-bound-only entries
(Remark~\ref{rem:288}, which corrects a misreading an earlier draft of our
results file contained). What they do not supply --- and what nobody supplies
--- is an artifact. The strongest lower bound previously \emph{reported} by a
method capable in principle of certifying is $d\ge11$, from Chen, Jafari and
Lai \cite{ChenJafariLai} at a $7200$\,s timeout, solver-asserted, with no
proof files in their repository. Our $\dX\ge14$ improves on that bound and is,
as far as the sweep of \S\ref{ssec:credit} could determine, the only
machine-checkable lower bound on record for this code. It falls four short of
$18$, and that shortfall is a limitation of our encoding, not evidence.

\smallskip
\noindent\textbf{(C) The rest of the family.} The values $d=6,10,10$ for
$[[72,12,6]]$, $[[90,8,10]]$, $[[108,8,10]]$, and the distances of Steane,
Golay, the rotated surface codes and the five-qubit code, are all prior art;
several are computed exactly, without certificates, in \cite{ChenJafariLai}.
For Steane, Golay, $[[72,12,6]]$ and $[[90,8,10]]$, replayable proof artifacts
also exist in the LEAN-QEC repository, so ``certified'' is shared ground
there and our contribution for those four is the trusted base, not primacy.
We include them because a certificate format is worth nothing without a
regression suite, and because six of them admit an independent brute-force
cross-check (\S\ref{ssec:brute}).

\subsection{Provenance and credit}\label{ssec:credit}
This subsection was written after a first-hand check of the prior art on
August 4, 2026, in the course of which two claims in an earlier internal draft
of these results turned out to be wrong. What follows is the corrected record.
Table~\ref{tab:credit} is the claim-by-claim version.

\emph{The codes.} The bivariate-bicycle family, including $\gross$, is Bravyi,
Cross, Gambetta, Maslov, Rall and Yoder's \cite{BCGMRY}. All five parity-check
matrices in our corpus were reconstructed from the construction in that paper
and are byte-identical to it (\S\ref{ssec:matrices}). Their distances were
computed there by the mixed-integer-programming method of Landahl, Anderson
and Rice \cite{LandahlAndersonRice}. Every distance value we certify for a BB
code, and every distance value we certify for Steane, Golay, the surface codes
and the five-qubit code, was already known. This note contributes
\emph{certificates for known numbers}, plus one number
($\dX\ge14$ for $n=288$) that is new only in the sense of being machine-checkable.

\emph{Machine-checked quantum distance proofs.} The approach is LEAN-QEC's
\cite{LeanQEC} (Ehatamm, Lee, Wu and Tao, 15 May 2026): a SAT encoding of the
distance problem whose refutation is replayed inside the Lean~4 kernel. Their
paper is candid about the gross code: it reports $\gross$ as dispatched by an
external \texttt{cvc5} call rather than in the kernel, states that the
$144$-qubit instance pushes their certificate-replay budget beyond what they
can verify in the kernel, and names kernel-side replay at that size as their
next concrete engineering target. \textbf{Their repository has since moved past
the paper.} At commit \texttt{c73827d} of \texttt{VerifiedQC/Lean-QEC}, dated
2026-07-10, a notes file reports that the $\gross$ BB code distance has been
fully verified with \texttt{bv\_decide} in about thirty minutes, by a pipeline
whose final step translates the LRAT proof of unsatisfiability back into Lean
and checks it with the kernel. We therefore make \emph{no} priority claim for a
machine-checked gross-code distance, and we have removed such a claim wherever
it appeared in our own working notes.

Several differences remain, and we state them as observations of a moving
repository at a single commit, which may already be stale. At
\texttt{c73827d}, the file \texttt{BB144.lean} contains two \texttt{sorry}
placeholders --- \texttt{BB144\_X\_ker\_rank} at line~$69$ and
\texttt{BB144\_Z\_ker\_rank} at line~$72$ --- through which the end-to-end
theorem \texttt{BB144\_dist\_12} routes, so that theorem is not
\texttt{sorry}-free as committed; three further lemmas use
\texttt{native\_decide}, which their own paper notes extends the trusted base
with Lean's compiler; and no LRAT artifact is committed for the $144$-qubit
instance, though LRATs are committed for smaller codes. Their $144$-qubit
encoding uses sorted-location symmetry constraints, i.e.\ it is
symmetry-broken; we found no counterpart anywhere to the symmetry-free proofs
of \S\ref{ssec:gross}. We note also, since it would be easy to overstate their
coverage in the other direction, that \texttt{BB108.lean} at the same commit
carries \texttt{sorry} at lines~$120$ and~$132$ with its
\texttt{bv\_decide} invocation commented out, so their kernel-checked ladder
should not be described as reaching $n=108$ either. None of this is a
criticism of work in progress; a repository ahead of its paper is a good
problem to have. It is the reason we describe our contribution as a trusted
base rather than a first, and we have made these statements as specific as we
can so that they can be checked and, when they go stale, corrected.

\emph{Mathematical programming, independently.} Cruz-Benito, Cross, Kremer and
Faro \cite{CruzBenito} (1 June 2026, IBM) revisit the gross-code distance with
a mixed-integer linear program and close it to a MIP gap of zero, confirming
$d=12$ exactly. That is a stronger form of the same kind of evidence as
\cite{BCGMRY}: an exhaustive search whose exhaustiveness is attested by the
solver's own optimality certificate rather than exported as a checkable
object. We cite it as the current reference point for the value, and note that
the gap it leaves --- a MILP optimality proof is not a file --- is precisely
the one this note fills.

\emph{SAT-based distance computation at scale.} Chen, Jafari and Lai
\cite{ChenJafariLai} (29 May 2026) compute BB distances with a battery of
solvers, obtaining $d=12$ for $\gross$ exactly and, for $[[288,12,18]]$, lower
bounds of $11$ under a $7200$\,s budget. Their results are solver-asserted ---
their repository \texttt{guluchen/QDistSAT} contains no proof or certificate
files --- and their $d\ge11$ is the strongest quantity previously published for
that code \emph{as a lower bound}, as against the exact value $d=18$ asserted
in \cite{BCGMRY}. Our $\dX\ge14$ improves on it and, for the first time,
backs a bound for this code with a replayable artifact.

\emph{Infrastructure.} The LRAT format and the checking discipline are due to
Cruz-Filipe, Heule, Hunt, Kaufmann and Schneider-Kamp \cite{LRAT}, building on
\texttt{drat-trim} \cite{DRATtrim}; our internal checker implements
RUP-with-hints against that format. The solver is CaDiCaL \cite{CaDiCaL}. The
cardinality encoding is Sinz's sequential counter \cite{Sinz05}; the XOR gates
are Tseitin's \cite{Tseitin68}. The CSS structure is Calderbank--Shor
\cite{CS96} and Steane \cite{Steane96}; the stabilizer formalism is
Gottesman's \cite{Gottesman97}.

\emph{Adjacent work on certificate import.} PBLean \cite{PBLean} imports
VeriPB pseudo-Boolean certificates into Lean~4. It is not quantum, and it
requires Lean, so it sits on the opposite side of design constraint (ii); but
it is the closest prior art we found to the general project of making
solver output externally checkable in a small trusted base, and we cite it as
such.

\begin{table}[t]
\small
\begin{tabular}{L{0.30\textwidth}L{0.40\textwidth}L{0.20\textwidth}}
\toprule
Item & Earliest source we verified & Status here \\
\midrule
The BB codes; $d=12$ for $\gross$; $d=18$ for $[[288,12,18]]$
& Bravyi et al.\ \cite{BCGMRY}, Table 3 (MIP method of
  \cite{LandahlAndersonRice})
& Re-verified, not claimed \\
$d=12$ for $\gross$ confirmed exactly by MILP, MIP gap $0$
& Cruz-Benito, Cross, Kremer, Faro \cite{CruzBenito}, 1 June 2026
& Re-verified, not claimed; no certificate emitted there \\
$\dX=\dZ$ for BB codes
& Bravyi et al.\ \cite{BCGMRY}, supplemental lemma
& Cited; only the explicit permutation certificate is ours \\
Kernel-checked SAT distance proofs for quantum codes
& LEAN-QEC \cite{LeanQEC}, 15 May 2026
& Not claimed \\
Machine-checked $\gross$ distance
& LEAN-QEC repository, commit \texttt{c73827d}, 2026-07-10
& Not claimed \\
$d$ computed exactly by SAT for $n\le144$ BB codes
& Chen--Jafari--Lai \cite{ChenJafariLai}, 29 May 2026
& Re-verified, not claimed \\
$d\ge11$ for $[[288,12,18]]$ by SAT
& Chen--Jafari--Lai \cite{ChenJafariLai}
& Improved to $\dX\ge14$, with artifact \\
\midrule
\multicolumn{3}{l}{\emph{Offered as new:}}\\
\multicolumn{3}{L{0.93\textwidth}}{The certificate format of
\S\ref{sec:formats} and its machine-checked exactness conditions; the
symmetry-free $\gross$ proofs; the only machine-checkable lower bound on
record for $[[288,12,18]]$, at $\dX\ge14$; the ZX-duality permutation ---
Bravyi et al.'s lemma, not ours --- packaged as a $\sim$$15$\,ms checkable
certificate; and a trusted base of $649$ standard-library Python lines with
no solver and no proof assistant in it.}\\
\bottomrule
\end{tabular}
\medskip
\caption{Claim-by-claim provenance.}
\label{tab:credit}
\end{table}

\subsection{What is not claimed}\label{ssec:notclaimed}
We do not claim any new distance \emph{value}. We do not claim priority for
machine-checked quantum distance proofs. We do not claim that
$[[288,12,18]]$ has $d<18$; our $\dX\ge14$ is a lower bound, four short of the
literature value, and the gap is a limitation of our encoding, not evidence.
We do not certify $k$: the code dimensions in Table~\ref{tab:main} are
recomputed, not certified, though side condition (c) of
Theorem~\ref{thm:exact} pins the logical dimension implicitly, and the audit
recomputed every $k$ independently. We do not claim the corpus is free of
defects: four are reported in \S\ref{sec:defects} --- one fixed in this
release, three not --- including a genuine latent soundness hole in a checker
branch that no shipped certificate exercises. And we do not claim novelty for
the framing itself: this is a verification contribution, and the reason to
read it is the audit of \S\ref{sec:audit}, not the numbers in
Table~\ref{tab:main}.

\section{Codes, distances, and the two certificate problems}\label{sec:prelim}

We work over $\FF$. A binary vector $x\in\FF^{n}$ has weight $\wt(x)$ and
support $\supp(x)$. For a matrix $M$, $\rowsp(M)$ is its row space, and for a
permutation $\pi$ of $\{0,\dots,n-1\}$ we write $M\pi$ for the
column-permuted matrix $(M\pi)_{r,i}=M_{r,\pi(i)}$.

\begin{definition}[CSS code]
A CSS code is a pair $\HX\in\FF^{m_{X}\times n}$, $\HZ\in\FF^{m_{Z}\times n}$
with $\HX\HZ^{\mathsf T}=0$. Its dimension is $k=n-\rank\HX-\rank\HZ$. An
\emph{$X$-logical} is a vector $x$ with $\HZ x=0$; it is \emph{nontrivial} if
$x\notin\rowsp(\HX)$. Set
\[
  \dX=\min\{\wt(x): \HZ x=0,\ x\notin\rowsp(\HX)\},
\]
and $\dZ$ symmetrically with $X$ and $Z$ interchanged. The code distance is
$d=\min(\dX,\dZ)$.
\end{definition}

The identity $d=\min(\dX,\dZ)$ is the standard CSS fact and we use it without
further comment; it is item~5 in the trusted-base inventory of
\S\ref{sec:tcb}. For a general (non-CSS) stabilizer code we use the symplectic
form $\omega$ on $\FF^{2n}$, $\omega(p,q)=\sum_{i}(p_{i}q_{n+i}+p_{n+i}q_{i})$,
a stabilizer matrix $S\in\FF^{m\times 2n}$ with pairwise-commuting rows, and
qubit weight $\wt_{q}(e)=\#\{i: e_{i}=1 \text{ or } e_{n+i}=1\}$.

Certifying $d$ splits into two problems of completely different character.
An \emph{upper} bound is existential and needs only a witness. A \emph{lower}
bound is universal: it asserts that a search space of size
$\binom{n}{\le K}$ intersected with a coset is empty, and no small object
witnesses that directly. The design question this note answers is what the
smallest honest replayable object for the second problem looks like.

\section{The certificate formats}\label{sec:formats}

\subsection{Upper bounds: witness pairs}\label{ssec:upper}
An $X$-sector upper-bound certificate is a pair $(x,z)\in\FF^{n}\times\FF^{n}$
together with a declared weight $w$. The checker
(\texttt{check\_witness.py}, $95$ lines) verifies
\[
  \HZ x = 0,\qquad \HX z = 0,\qquad x\cdot z = 1,\qquad \wt(x)=w .
\]
Soundness is immediate: $\HZ x=0$ makes $x$ an $X$-logical; every element of
$\rowsp(\HX)$ is orthogonal to $z$ because $\HX z=0$, so $x\cdot z=1$ forces
$x\notin\rowsp(\HX)$. Hence $\dX\le w$. Three matrix--vector products; no
solver, no proof, no search. The gross-code witnesses are $970$-byte files and
check in $7$\,ms.

The general stabilizer form is the same shape with $\omega$ in place of the
dot product: $(e,f)$ with $\omega(s,e)=\omega(s,f)=0$ for every stabilizer row
$s$, and $\omega(e,f)=1$, certifies $d\le\wt_{q}(e)$. The five-qubit code
carries this variant end to end.

\subsection{Lower bounds: a CNF the checker writes itself}\label{ssec:lower}
Fix a sector, say $X$, and a bound $K$. The canonical CNF
$\Phi_{K}(\HZ,P)$ over variables $x_{1},\dots,x_{n}$ (plus auxiliaries)
asserts
\begin{enumerate}
\item[(1)] $\HZ x=0$, one Tseitin XOR chain per row of $\HZ$, with the chain
output forced false;
\item[(2)] $b_{j}=z_{j}\cdot x$ for each row $z_{j}$ of a \emph{pairing}
matrix $P$, one XOR chain each, output free;
\item[(3)] the clause $(b_{1}\vee\cdots\vee b_{p})$;
\item[(4)] a Sinz sequential at-most-$K$ counter \cite{Sinz05} on
$x_{1},\dots,x_{n}$;
\item[(5)] optionally, forced unit literals (\S\ref{ssec:sym}).
\end{enumerate}
A lower-bound certificate is a JSON file naming $\HX,\HZ,P$, the value $K$,
the sector, and for each instance a list of forced literals and an LRAT file.
The checker \texttt{check\_lower.py} ($481$ lines) does four things: verifies
the side conditions of Theorem~\ref{thm:exact}; \emph{regenerates}
$\Phi_{K}(\HZ,P)$ from the raw matrices in a fixed variable order; replays the
LRAT against its own clause list; and, for symmetry-broken certificates,
verifies the hypotheses of Lemma~\ref{lem:orbit}.

The shipped \texttt{.cnf} files are never read. They are informational, and the
audit confirmed this by appending a contradictory clause pair to one of them
and observing the check pass unchanged (\S\ref{sec:audit}).

\begin{theorem}[Exactness of the CSS encoding]\label{thm:exact}
Let $\HX\in\FF^{m_{X}\times n}$, $\HZ\in\FF^{m_{Z}\times n}$ and
$P\in\FF^{p\times n}$ satisfy
\begin{enumerate}
\item[(a)] $\HX\HZ^{\mathsf T}=0$;
\item[(b)] $\HX P^{\mathsf T}=0$;
\item[(c)] $\rank\begin{pmatrix}\HZ\\P\end{pmatrix}=n-\rank\HX$.
\end{enumerate}
Then
\[
  \{x\in\FF^{n}: \HZ x=0,\ Px\ne 0\}
  \;=\;\ker\HZ\setminus\rowsp(\HX),
\]
the set of nontrivial $X$-logicals. Consequently $\Phi_{K}(\HZ,P)$ is
satisfiable if and only if there is a nontrivial $X$-logical of weight at most
$K$; an unsatisfiability proof for $\Phi_{K}(\HZ,P)$ certifies $\dX\ge K+1$.
\end{theorem}

\begin{proof}
Put $V=\ker\HZ\cap(\rowsp P)^{\perp}=(\rowsp\HZ+\rowsp P)^{\perp}$.
By (a), $\rowsp(\HX)\subseteq\ker\HZ$; by (b), $\rowsp(\HX)\subseteq
(\rowsp P)^{\perp}$. Hence $\rowsp(\HX)\subseteq V$. By (c),
\[
  \dim V \;=\; n-\rank\begin{pmatrix}\HZ\\P\end{pmatrix}\;=\;\rank\HX
  \;=\;\dim\rowsp(\HX),
\]
so $V=\rowsp(\HX)$. Now for $x\in\ker\HZ$ we have $Px\neq 0$ iff
$x\notin(\rowsp P)^{\perp}$ iff $x\notin V=\rowsp(\HX)$, which is the claim.
The second sentence follows because clauses (1)--(3) of $\Phi_{K}$ encode
exactly $\HZ x=0\wedge Px\neq0$ and clause group (4) encodes $\wt(x)\le K$.
\end{proof}

Both directions of Theorem~\ref{thm:exact} are load-bearing and for different
reasons. Condition (b) gives \emph{soundness}: without it a satisfying
assignment need not be a nontrivial logical, and UNSAT would prove nothing
about the code. Condition (c) gives \emph{completeness}: without it the
encoding could miss logicals, and UNSAT would prove less than claimed. The
checker verifies (a) and (b) by explicit $\FF$ inner products and (c) by two
Gaussian eliminations over bitmask integers, on every certificate, before it
looks at the proof. This is why the pairing matrix $P$ --- which the pipeline
produced, and which a hostile pipeline could try to weaken --- is not trusted
input. Negative control 7 of \S\ref{sec:audit} deletes one of its twelve rows
and the checker rejects on condition (c).

\subsection{Symmetry breaking, and the orbit lemma}\label{ssec:sym}
For the largest instances the search space can be cut by the code's own
automorphisms. A symmetry-broken certificate ships permutations
$\pi_{1},\dots,\pi_{r}$, a block boundary $b$, and \emph{two} instances with
prescribed forced literals.

\begin{lemma}[Orbit lemma]\label{lem:orbit}
Let $G\le S_{n}$ be generated by permutations $\pi$ such that
\begin{enumerate}
\item[(i)] $\pi$ preserves the partition $\{0,\dots,b-1\}\sqcup\{b,\dots,n-1\}$;
\item[(ii)] $\rowsp(\HX\pi)=\rowsp(\HX)$ and $\rowsp(\HZ\pi)=\rowsp(\HZ)$;
\end{enumerate}
and suppose $G$ acts transitively on each of the two blocks. If a nontrivial
$X$-logical of weight $\le K$ exists, then one exists whose support either
contains qubit $0$, or is disjoint from $\{0,\dots,b-1\}$ and contains qubit
$b$.
\end{lemma}

\begin{proof}
By (ii), $x\mapsto x\circ\pi$ maps $\ker\HZ$ to itself and $\rowsp(\HX)$ to
itself, hence maps nontrivial $X$-logicals to nontrivial $X$-logicals, and it
preserves weight. Let $x$ be nontrivial of weight $\le K$. If
$\supp(x)$ meets the left block, pick $i$ in the intersection and $g\in G$
with $g(0)=i$, possible by transitivity; then $y=x\circ g$ has $y_{0}=1$.
Otherwise $\supp(x)$ meets the right block, say at $i\ge b$; pick $g\in G$
with $g(b)=i$; then $y=x\circ g$ has $y_{b}=1$, and $y$ still vanishes on the
left block because $g$ preserves blocks.
\end{proof}

The two cases are exactly the two forced-literal patterns the checker demands:
instance $0$ has \texttt{forced} $=[\,1\,]$ and instance $1$ has
\texttt{forced} $=[-1,\dots,-b,\,b+1]$. Everything Lemma~\ref{lem:orbit}
hypothesizes is machine-checked: that each shipped array is a permutation, that
it preserves the blocks (index comparison), that it is a rowspace automorphism
of both $\HX$ and $\HZ$ (rank identity $\rank(H)=\rank(H\cup H\pi)$),
that the generated group is block-transitive (breadth-first search over the
generators and their inverses), and that the two forced patterns are literally
the two above. Only the four-line implication in the proof is human-verified.
For $\gross$ and $[[288,12,18]]$ the two shipped permutations were
independently identified in the audit as the torus translations $x^{1}y^{0}$
and $x^{0}y^{1}$.

\begin{remark}\label{rem:symfree}
The gross code does not need this lemma. The $868$\,MB $X$ proof and the
$672$\,MB $Z$ proof are single-instance certificates with empty
\texttt{forced} lists, and the checker's non-symmetry branch asserts exactly
that: one instance, no forced literals. So $d(\gross)=12$ is certified with
the orbit lemma removed from the trusted base. The symmetry-broken $X$
certificate is retained because it is $7\times$ smaller and $17\times$ faster
to replay, and because it demonstrates the mechanism at $n=144$ before it is
used in earnest at $n=288$.
\end{remark}

\subsection{The ZX-duality certificate}\label{ssec:duality}
For BB codes $\dX=\dZ$. This is not our observation: it is a lemma in the
supplemental material of Bravyi et al.\ \cite{BCGMRY}, invoked there in
exactly the way we invoke it --- to halve the work --- with the remark that
the hypothesis in question ``is satisfied for BB LDPC codes''. What we add is
an artifact: an explicit permutation, shipped as a file, that turns the lemma
into something a reader checks in about $15$\,ms instead of something a reader
believes.

\begin{lemma}[Duality]\label{lem:dual}
Let $\Pi$ be a permutation of the $n$ qubits with
$\rowsp(\HX\Pi)=\rowsp(\HZ)$ and $\rowsp(\HZ\Pi)=\rowsp(\HX)$. Then
$x\mapsto y$, $y_{i}=x_{\Pi(i)}$, is a weight-preserving bijection from
nontrivial $X$-logicals to nontrivial $Z$-logicals, and $\dX=\dZ$.
\end{lemma}

\begin{proof}
Applying $\Pi^{-1}$ to the two hypotheses gives
$\rowsp(\HX)=\rowsp(\HZ\Pi^{-1})$ and $\rowsp(\HZ)=\rowsp(\HX\Pi^{-1})$.
Let $h$ be a row of $\HX$. Then $h\cdot y=\sum_{i}h_{i}x_{\Pi(i)}
=\sum_{j}h_{\Pi^{-1}(j)}x_{j}=(h\circ\Pi^{-1})\cdot x$, and
$h\circ\Pi^{-1}\in\rowsp(\HX\Pi^{-1})=\rowsp(\HZ)$, so $h\cdot y=0$ whenever
$\HZ x=0$. Thus $\HX y=0$. If $y\in\rowsp(\HZ)$, say $y=c\HZ$, then
$x_{j}=y_{\Pi^{-1}(j)}$ exhibits $x\in\rowsp(\HZ\Pi^{-1})=\rowsp(\HX)$,
contradicting nontriviality of $x$. Weight is preserved because $\Pi$ permutes
coordinates. The same argument with $X$ and $Z$ interchanged gives the inverse
map, so the correspondence is a bijection and $\dX=\dZ$.
\end{proof}

The certificate is the permutation $\Pi$; the checker
(\texttt{check\_duality.py}, $73$ lines) does two rank comparisons. The corpus
ships one for each of the five BB codes, and each checks in a few milliseconds
to a few tens of milliseconds --- $3.8$, $4.7$, $6.1$, $9.9$ and $37.9$\,ms
for $n=72,90,108,144,288$ on a loaded machine. In every case $\Pi$ is the
composite of the block swap with $(a,b)\mapsto(-a,-b)$ on the torus indices.

\begin{remark}[The $n=288$ certificate, and when it appeared]\label{rem:no288dual}
The audit of \S\ref{sec:audit} found no duality certificate for
$[[288,12,18]]$ (discrepancy D2 of \S\ref{sec:defects}): the generator listed
the code but had not been re-run, so what the corpus certified there was
$14\le\dX\le18$ rather than $14\le d\le18$. The auditor independently
confirmed that the same permutation family does satisfy both rowspace
identities at $n=288$, so the mathematics was never in doubt --- only the
artifact was missing, which is exactly the condition this paper exists to
avoid. The generator has since been re-run and
\texttt{bb288/duality.json} now ships. It is the one certificate in the
release that is \emph{not} among the audit's $47$ checks and not among the
$182$ hashes in \texttt{manifest.json}, both of which predate it; it passes
\texttt{check\_duality.py}, and its permutation was verified a second time,
from scratch, against independently loaded matrices. With it,
$14\le d\le 18$ is certified rather than merely true, and
Table~\ref{tab:main} is worded accordingly. No $Z$-sector witness is needed:
$d=\min(\dX,\dZ)=\dX$ once duality is in hand.
\end{remark}

\subsection{The symplectic variant}\label{ssec:symp}
For a general stabilizer code the encoding uses $3n$ variables --- $u_{i}$,
$v_{i}$ for the $X$- and $Z$-components and $q_{i}$ for ``qubit $i$ is in the
support'', tied by $(\lnot u_{i}\vee q_{i})$, $(\lnot v_{i}\vee q_{i})$,
$(u_{i}\vee v_{i}\vee\lnot q_{i})$ --- with XOR chains for the symplectic
products and the at-most-$K$ counter on the $q$'s.

\begin{proposition}\label{prop:symp}
Let $S\in\FF^{m\times 2n}$ have pairwise $\omega$-orthogonal rows, let
$L\in\FF^{\ell\times2n}$ satisfy $\omega(s,\lambda)=0$ for all rows
$s$ of $S$ and $\lambda$ of $L$, and suppose
$\rank\binom{S}{L}=\rank S+\ell=n+k$ where $k=n-\rank S$. Then
$\{e:\omega(s,e)=0\ \forall s,\ \exists\lambda\ \omega(e,\lambda)=1\}$ is
exactly the set of nontrivial logical operators.
\end{proposition}

\begin{proof}
Let $C$ be the $\omega$-centralizer of $\rowsp(S)$, of dimension
$2n-\rank S=n+k$. The radical of $\omega|_{C}$ is $\rowsp(S)$. The rows of $S$
and $L$ all lie in $C$ and span a subspace of dimension $\rank S+\ell=n+k$, so
they span $C$. If $e\in C$ and $\omega(e,\lambda)=0$ for every row of $L$,
then $\omega(e,\cdot)$ vanishes on $C$, so $e$ lies in the radical, i.e.\
$e\in\rowsp(S)$. The converse inclusion is immediate.
\end{proof}

The five-qubit code $[[5,1,3]]$ carries this path end to end: a $128$-byte
witness and a $1{,}514$-byte LRAT for $d\ge3$. It exists to keep the non-CSS
branch of the checker exercised; see also defect D4 in \S\ref{sec:defects},
which is a missing guard in exactly that branch.

\section{The trusted base}\label{sec:tcb}

A certificate is only as good as the list of things its reader must still
believe. Here is that list, split into what is re-verified per certificate by
code one can read (\textsc{checked}) and what is argued in prose and verified
nowhere (\textsc{assumed}).

\subsection{Software}
Three files, $649$ lines of Python in total:
\texttt{check\_witness.py} ($95$), \texttt{check\_duality.py} ($73$),
\texttt{check\_lower.py} ($481$). Their only imports are
\texttt{gzip}, \texttt{json}, \texttt{os}, \texttt{subprocess}, \texttt{sys},
\texttt{tempfile} and \texttt{time}. No \texttt{numpy}, no
compiled helper, no network. CPython and the operating system must be trusted.
If the optional \texttt{--external} path is used to delegate LRAT replay to a
compiled checker, that binary and the temporary-file marshalling re-enter the
trusted base; the audit of \S\ref{sec:audit} never used it, so for that audit
the trusted base contained no compiled code at all.

\emph{Not} trusted, and demonstrably so: CaDiCaL; the generating pipeline
(\texttt{certify.py}, \texttt{qec\_lib.py}, \texttt{gen\_duality.py},
\texttt{manifest.py}); the shipped \texttt{.cnf} files; \texttt{meta.json};
\texttt{manifest.json}; and the prose of this paper. All of them can be
deleted and every certificate still verifies.

\begin{remark}[The checker moved after the audit]\label{rem:checkerdrift}
The auditor of \S\ref{sec:audit} read a $419$-line
\texttt{check\_lower.py}; the file shipped with this release is $481$ lines.
The difference is one addition: an optional truncated Bailleux--Boufkhad
totalizer as an alternative to the Sinz cardinality encoding, put in for
larger $[[288,12,18]]$ rungs that are still running and are not part of this
release. It is selected only by a certificate that declares
\texttt{"cardinality": "totalizer"}, and \emph{no certificate in this release
declares it}: all take the Sinz default, which is the code path the audit
read and exercised. A reader minimising trusted-base surface can delete the
totalizer branch and re-run everything. We flag the drift rather than quietly
re-using the audit's line count, because the size of that number is one of the
claims. To close the loop, the entire corpus was re-run against the shipped
$481$-line checker before release --- $20$ witnesses, $5$ duality certificates
and $23$ lower-bound certificates, $48$ in all, every one accepted, including a
second pure-Python replay of the $2.94$\,GB proof. Those re-runs were made on a
machine under heavy load and their wall-clock times are correspondingly three
to five times the audit's; the pass/fail outcomes are what they establish.
\end{remark}

\subsection{Machine-checked, per certificate}
\begin{itemize}
\item CSS orthogonality $\HX\HZ^{\mathsf T}=0$ --- condition (a).
\item Every pairing row in the kernel of the quotient matrix --- condition (b),
soundness.
\item The rank/spanning identity --- condition (c), completeness.
\item For the symplectic branch: commutativity of $S$, $L$ in the centralizer,
and the two rank identities of Proposition~\ref{prop:symp}.
\item For \texttt{\_sym} certificates: that each permutation is a permutation,
preserves blocks, is a rowspace automorphism of both $\HX$ and $\HZ$; that the
generated group is block-transitive; and that the forced patterns are exactly
the two Lemma~\ref{lem:orbit} licenses.
\item For duality certificates: both rowspace identities.
\item Regeneration of the entire CNF from the raw matrices, and RUP replay of
the LRAT against \emph{that} clause list.
\end{itemize}

\subsection{Assumed}
\begin{enumerate}
\item \emph{The Sinz sequential at-most-$K$ counter is complete}
\cite{Sinz05}: every weight-$\le K$ assignment extends to a satisfying
assignment of the counter variables. If it were accidentally
over-constraining, UNSAT would mean less than claimed. The emitted clause set
was read against \cite{Sinz05} and is the standard one.
\item \emph{The Tseitin XOR gate is definitional} \cite{Tseitin68}: the four
clauses per gate encode $u=a\oplus b$ and are always extendable.
\item \emph{The orbit lemma}, Lemma~\ref{lem:orbit} --- four lines, proved
above, machine-checked hypotheses, human-verified implication. Used by the
symmetry-broken $\gross$ $X$ certificate and all three $[[288,12,18]]$ rungs.
\textbf{Not used by} the symmetry-free $\gross$ certificates
(Remark~\ref{rem:symfree}).
\item \emph{The duality lemma}, Lemma~\ref{lem:dual} --- likewise, checked
hypotheses, human-verified implication.
\item \emph{The CSS fact $d=\min(\dX,\dZ)$}, used to turn a certified pair
into a statement about the code. No script touches it.
\item \emph{LRAT semantics.} The internal checker implements RUP-with-hints:
negate the lemma, unit-propagate through the hinted clauses in order, demand a
conflict. It skips hints that are already satisfied or non-unit --- sound, since
skipping can only fail to find a conflict, never invent one --- and it
\emph{refuses} negative (RAT) hints outright, which incidentally establishes
that every proof in the corpus is pure RUP. It requires an empty clause to be
derived and verified.
\item \emph{That the matrices are the intended code.} No certificate
establishes this; the checkers take \texttt{HX.txt} and \texttt{HZ.txt} at face
value. Section~\ref{ssec:matrices} closes this for the five BB codes at byte
level.
\end{enumerate}

\noindent In one sentence: a skeptic must believe that three short
standard-library Python files do what they appear to do, that CPython and the
OS are not lying, that the Sinz and Tseitin encodings are standard and
complete, that the four-line duality lemma and the four-line orbit lemma are
correct, and that $d=\min(\dX,\dZ)$. Everything else can be thrown away.

\section{The gross code}\label{ssec:gross}

$\gross$ is the $\ell=12$, $m=6$ member of the BB family, with
$A=x^{3}+y+y^{2}$ and $B=y^{3}+x+x^{2}$, $\HX=[A\mid B]$,
$\HZ=[B^{\mathsf T}\mid A^{\mathsf T}]$ \cite{BCGMRY}. Its certified distance
decomposes as follows; all byte counts are the shipped files.

\begin{center}
\small
\begin{tabular}{llrrr}
\toprule
statement & artifact & bytes & solver (s) & replay \\
\midrule
$\dX\le 12$ & \texttt{witness\_X.json} & $970$ & $0.025$ & $7.4$\,ms\\
$\dZ\le 12$ & \texttt{witness\_Z.json} & $970$ & $0.029$ & $7.3$\,ms\\
$\dX\ge 12$, symmetry-broken & \texttt{lower\_X\_K11\_sym.json} $+2$ LRATs
  & $123{,}857{,}070$ & $45.0$ & $10.1$\,s\\
$\dZ\ge 12$, symmetry-free & \texttt{lower\_Z\_K11.json} $+1$ LRAT
  & $671{,}988{,}205$ & $227.3$ & $72.9$\,s\\
$\dX\ge 12$, symmetry-free & \texttt{lower\_X\_K11.json} $+1$ LRAT
  & $867{,}803{,}294$ & $342.3$ & $176.4$\,s\\
$\dX=\dZ$ & \texttt{duality.json} & $103+466$ & --- & $9.9$\,ms\\
\bottomrule
\end{tabular}
\end{center}

\noindent Total solver time for the complete symmetry-free route is
$569.6$\,s: under ten minutes on a laptop. Total replay time for that route,
in pure Python with no compiled code anywhere, is $249.3$\,s. The two
symmetry-free proofs are what let us say that $d(\gross)=12$ holds on the
trusted base of \S\ref{sec:tcb} \emph{minus} item~3.

The value $d=12$ is not new; it is \cite[Table 3]{BCGMRY}, obtained by mixed
integer programming \cite{LandahlAndersonRice}, confirmed exactly at MIP gap
zero by \cite{CruzBenito}, and reproduced by SAT in
\cite{ChenJafariLai}. A machine-checked proof of it is not new either: the
LEAN-QEC repository reports one as of 2026-07-10 (\S\ref{ssec:credit}). What is
offered here is the shape of the artifact --- an $868$\,MB file plus a
$481$-line reader, with no symmetry hypothesis, no solver and no proof
assistant --- and the fact that it replays inside $51$\,MB of memory.

\begin{remark}[Timing variability]
The replay times above are the audit's measurements on an otherwise idle
machine. A re-run of the symmetry-broken $X$ certificate performed while
writing this note, under load, took $30.4$\,s rather than $10.1$\,s. All timing
figures in this paper should be read with a factor of three of slack; the byte
counts and the pass/fail outcomes are exact.
\end{remark}

\section{$[[288,12,18]]$}\label{sec:288}

For the $\ell=m=12$ member, $A=x^{3}+y^{2}+y^{7}$, $B=y^{3}+x+x^{2}$, we
certify a ladder in the $X$ sector, each rung an independently valid
symmetry-broken certificate:
\[
  \dX\ge 10\ (65\,\text{MB},\ 11.0\,\text{s solver}),\quad
  \dX\ge 12\ (351\,\text{MB},\ 85.6\,\text{s}),\quad
  \dX\ge 14\ (2.94\,\text{GB},\ 512.6\,\text{s}),
\]
together with a weight-$18$ witness for $\dX\le18$ and a duality certificate
(Remark~\ref{rem:no288dual}) that converts both into statements about $d$. The
$K=13$ proof is $2{,}941{,}958{,}076 + 191{,}479$ bytes across the two orbit
instances and replays in pure Python in $413.6$\,s with a peak resident set of
$79$\,MB --- the memory figure being the interesting one, since the proof is
more than thirty times larger than the memory used to check it.

\begin{remark}[What the literature says, stated correctly]\label{rem:288}
An earlier internal draft of our results claimed that Bravyi et al.\ screened
this code with a heuristic and conceded that its $d\le18$ ``is unlikely to be
tight''. That is a misreading and we correct it here. In \cite{BCGMRY}, the
quoted caveat concerns the \emph{circuit-level} distance $\dcirc\le18$, a
different quantity from the code distance. Table~3 of that paper lists
$[[288,12,18]]$ with no ``$\le$'', while $[[360,12,\le24]]$ and
$[[756,16,\le34]]$ carry one, and the caption states that the notation
``$\le d$'' marks entries for which only an upper bound is known; the
supplemental material states that the actual distance ``of each candidate code
was computed using the integer linear programming method''. So \cite{BCGMRY}
asserts $d=18$ exactly for this code, by ILP and without a checkable artifact.
Our interval $[14,18]$ is therefore \emph{not} new information about the
value, and any suggestion that it is the first distance information produced
for this code --- a suggestion the earlier draft made --- is withdrawn. What
$\dX\ge14$ is, is the strongest lower bound for this code that anyone can
check.
\end{remark}

The strongest quantity previously published for this code as a \emph{lower
bound} is $d\ge11$, from Chen, Jafari and Lai \cite{ChenJafariLai}, obtained
by several solvers under $7200$\,s timeouts and asserted rather than
certified; their public repository \texttt{guluchen/QDistSAT} contains no
proof artifacts. Our $\dX\ge14$ improves on it and adds what neither it nor
\cite{BCGMRY} has: a file. As far as the sweep recorded in
\S\ref{ssec:credit} could determine, it is the only machine-checkable lower
bound on record for $[[288,12,18]]$ at any strength.

Where the wall is, concretely: proof size grows roughly $8\times$ per ladder
rung on this encoding, so $K=15$ is an overnight job at an estimated
$25$\,GB, and $K=17$ --- the rung that would close the gap to $18$ --- is out
of reach of this encoding on this machine. Better encodings, in particular
the location-indexed encoding of \cite{LeanQEC}, proof compression, and
per-rung parallelism are the obvious next moves, and none of them is ours.

\section{Independent verification}\label{sec:audit}

The corpus was re-checked by a separate agent instance with no access to the
pipeline and no shared code, on the same machine, using
\texttt{/usr/bin/python3} (CPython 3.9.6). Its report is
\texttt{INDEPENDENT-VERIFICATION.md}. Summary: $47$ checks run, $47$ passed,
$0$ failed; $5{,}459{,}315{,}046$ bytes ($5.08$\,GiB) of LRAT replayed in
approximately $746$\,s of wall clock; peak resident set $79$\,MB; all $182$
SHA-256 entries in \texttt{manifest.json} re-hashed and matching. Because
\texttt{lrat-check} is not present in the repository (defect D3), \emph{every}
proof --- including the $2.94$\,GB one --- was replayed with the internal
pure-Python checker. That is the stronger result. Eleven negative controls,
built by the auditor on scratch copies, were all correctly rejected
(\S\ref{ssec:neg}), and six codes were cross-checked by brute force
(\S\ref{ssec:brute}). This audit, not the distance values, is the load-bearing
part of the paper.

\subsection{Are these the right matrices?}\label{ssec:matrices}
This is the check the certificates cannot perform on themselves, and it was
done first. The auditor rebuilt $\HX$ and $\HZ$ for all five BB codes directly
from the construction in \cite{BCGMRY} --- $n=2\ell m$, $x=S_{\ell}\otimes
I_{m}$, $y=I_{\ell}\otimes S_{m}$ over $\FF$ with $S_{r}$ the cyclic shift,
$\HX=[A\mid B]$, $\HZ=[B^{\mathsf T}\mid A^{\mathsf T}]$ --- using that
paper's Table~3 parameters. For every one of \texttt{bb72}, \texttt{bb90},
\texttt{bb108}, \texttt{bb144}, \texttt{bb288}, the shipped \texttt{HX.txt}
and \texttt{HZ.txt} are \textbf{byte-identical} to the reconstruction: no row
operations, no column permutation, no block swap. Independently recomputed
$k=n-\rank\HX-\rank\HZ$ gives $12,8,8,12,12$, matching the published
parameters, with row weight $6$ and column weight $3$ throughout. Every
algebraic side condition was then re-run a second time, in the auditor's own
implementation, against the \emph{reconstructed} matrices rather than the
shipped ones; all agreed.

\subsection{Brute force where brute force is possible}\label{ssec:brute}
For the six codes small enough, the auditor computed the minimum distance by
exhaustive Gray-code enumeration over a kernel basis, with code sharing
nothing with either the pipeline or the checkers: Steane, surface $d=3$,
surface $d=5$, Golay, the five-qubit code, and surface $d=7$ (a $2^{25}$
search, $242$\,s). Six exact distances, six agreements with the certified
values. From $[[72,12,6]]$ upward the search space is $2^{42}$ and beyond,
which is the entire point of the certificates.

\subsection{Negative controls}\label{ssec:neg}
Eleven tampered artifacts were built on scratch copies and all eleven were
rejected, each with the specific error one would want: a flipped witness bit
(``witness does not commute''), an understated weight, a relabelled $K$
(``LRAT lemma 2015 NOT verified''), a deleted empty clause, a truncated proof,
a live-but-wrong hint id, a gutted pairing matrix (rejected on condition (c),
``$31\ne42$''), a flipped bit of $\HX$ (rejected on condition (a)), a swapped
sector label, a corrupted hint id, and an identity duality permutation. In
addition, a contradictory clause pair was appended to
\texttt{bb72/lower\_X\_K5.cnf} and the check passed unchanged and in the same
time --- a positive control on the claim that the shipped CNF is genuinely
never read.

\section{Defects found}\label{sec:defects}

Reported verbatim, because a paper about trusted bases that suppresses its own
audit findings is not one.

\smallskip
\noindent\textbf{D1 (documentation).} The results file cites a
\texttt{tamper\_test/} directory that does not exist in the repository. The
claim is true --- the controls were rebuilt from scratch and all rejected
(\S\ref{ssec:neg}) --- but the citation is dangling.

\smallskip
\noindent\textbf{D2 (real, touched a headline, now closed).} At audit time
there was no duality certificate and no $Z$-witness for $[[288,12,18]]$, so
the corpus certified $14\le\dX\le18$, not $14\le d\le18$: the generator listed
that code and had evidently not been re-run. The underlying mathematics was
verified independently by the auditor, so the claim was true and uncertified,
which is precisely the condition this paper exists to avoid. This is the one
defect of the four that is fixed in this release: the generator has been
re-run, \texttt{bb288/duality.json} ships, it passes
\texttt{check\_duality.py}, and its permutation has been re-verified from
scratch. The caveats are that this certificate postdates the audit and is
therefore outside its $47$ checks, and that \texttt{manifest.json} --- whose
$182$ entries the audit re-hashed --- has deliberately not been regenerated,
so it does not list the new file. See Remark~\ref{rem:no288dual}.

\smallskip
\noindent\textbf{D3 (documentation).} \texttt{run\_all.sh} requires
\texttt{tools-drat-trim/lrat-check}, which is absent with no source or URL, so
the \texttt{lrat-check} timings previously quoted are not reproducible from
the repository as shipped. The pure-Python path is reproducible and is what
Table~\ref{tab:main} reports. Two earlier pure-Python figures were also stale
in the conservative direction ($21$\,s quoted versus $6.2$\,s measured for
$[[108,8,10]]$; $17.3$\,s versus $10.1$\,s for the symmetry-broken $\gross$
certificate), and the quoted throughput of $3.4$\,MB/s is closer to $12$\,MB/s.

\smallskip
\noindent\textbf{D4 (latent soundness hole).} In the CSS branch,
\texttt{check\_lower.py} requires a certificate with no symmetry block to have
exactly one instance with an empty forced list. The symplectic branch has
\emph{no such guard}: it passes the certificate's forced literals straight into
the encoder, which emits them as unit clauses. The auditor demonstrated this by
adding \texttt{"forced": [-1,-2,-3]} to the five-qubit certificate and
obtaining \texttt{OK lower bound: d >= 3} from a proof that now covers only
assignments avoiding those qubits. No shipped certificate is affected --- the
five-qubit certificate's forced list is empty, and this was checked --- but the
checker would accept a forged one. The one-line fix is to mirror the CSS
assertion. Two lesser observations in the same vein: the symmetry branch does
not assert that there are exactly two instances (harmless, since each is
replayed independently), and the CSS branch does not validate the sector label
(a typo silently selects the other branch, whose side conditions then fail, as
negative control 9 shows).

\smallskip
\noindent D4 is the interesting one. It is exactly the class of defect that the
certificate discipline is supposed to make findable: a hole in $481$ lines of
readable Python, found by reading them, rather than a hole in a solver.

\section{Relation to other work}\label{sec:related}

\emph{Heuristic upper bounds.} QDistRnd \cite{QDistRnd} produces upper bounds
with no performance guarantee; Stim's logical-error search \cite{Stim} is
documented as heuristic. Both are excellent for what they are. Our upper
bounds differ only in shipping the witness.

\emph{Exact methods without artifacts.} The mixed-integer-programming method
of Landahl, Anderson and Rice \cite{LandahlAndersonRice} is what
\cite{BCGMRY} used and remains the reference for the BB distances;
Cruz-Benito, Cross, Kremer and Faro \cite{CruzBenito} re-close the gross-code
MILP to gap zero. Chen, Jafari and Lai \cite{ChenJafariLai} push SAT and
MaxSAT to this family at scale and explicitly do not log proofs. Webster, Jacob and Higgott \cite{WJH}
survey the field and distinguish exact from heuristic methods carefully; their
desiderata do not include exportable certificates, and an exhaustive search of
their text turns up a single occurrence of certification, referring to
internal optimality rather than to an artifact. None of these produces
something a third party can replay. That gap, not the numbers, is what this
note addresses.

\emph{Proof-assistant approaches.} LEAN-QEC \cite{LeanQEC} is the closest work
and the one we have most to learn from: their location-indexed encoding is
better than ours, and we did not need it only because we stopped at $n=288$.
The design difference is where the trust sits. Their artifact is replayable
inside the Lean~4 kernel, which is a stronger guarantee than ours about the
\emph{checker}, at the cost of requiring Lean and Mathlib --- a toolchain fetch
--- to obtain it. Ours is replayable by any CPython, which is a weaker
guarantee about the checker and a much smaller barrier to a skeptic. Neither
dominates. PBLean \cite{PBLean} imports VeriPB certificates into Lean~4 and is
the nearest general prior art for the certificate-import problem.

\emph{What we could not find.} A search of the public record on 2026-08-04 ---
arXiv, code search for LRAT together with quantum distance, and the issue
trackers of Stim, \texttt{panqec}, \texttt{ldpc} and \texttt{qLDPC} --- turned
up no standalone quantum distance certificate requiring neither a solver nor a
proof assistant to check. A negative cannot be proved this way; we record the
scope of the search rather than a claim.

\section{Reproducibility}\label{sec:repro}

The artifacts sit beside this note in \texttt{qec/}. To re-check anything:

\begin{verbatim}
gunzip certificates/bb144/*.lrat.gz
python3 check_witness.py  certificates/bb144/witness_X.json
python3 check_lower.py    certificates/bb144/lower_X_K11.json
python3 check_lower.py    certificates/bb144/lower_Z_K11.json
python3 check_duality.py  certificates/bb144/duality.json
python3 check_lower.py    certificates/bb288/lower_X_K13_sym.json
\end{verbatim}

\noindent No virtual environment, no packages, no compiled binary. Any
CPython~3.8 or later should do; the audit used the macOS system interpreter,
3.9.6. Expect $176$\,s, $73$\,s and $414$\,s respectively for the three large
replays on an idle M4, more under load. \texttt{manifest.json} carries the
SHA-256 and byte count of all $182$ artifacts and
\texttt{manifest.py} re-checks them.

\smallskip
\noindent\emph{What the public repository can and cannot hold.} Four of the
proofs --- the two symmetry-free $\gross$ certificates and the first instance
of each of the $K=11$ and $K=13$ rungs at $n=288$ --- are between $79$\,MB and
$646$\,MB compressed and are not carried in git. Everything else is: the
descriptors, the parity-check matrices, the pairing and permutation files, the
witnesses, the duality certificates, the CNF inputs, the manifest, and every
proof small enough to ship, up to and including the $30$\,MB symmetry-broken
$\gross$ certificate and the $K=9$ rung at $n=288$. So a reader who clones and
runs nothing but CPython can still replay a certified $d=3$ for all four
sanity-tier codes, $d=7$ for Golay and the $d{=}7$ surface code, $d=6$ for
$[[72,12,6]]$, $d=10$ for $[[90,8,10]]$ and $[[108,8,10]]$, $d=12$ for the
gross code --- by the symmetry-broken $X$ certificate together with the
duality certificate and the weight-$12$ witness, the two symmetry-free proofs
being exactly the ones that do not fit --- and $\dX\ge10$ at $n=288$. For the
four omitted proofs the release ships \texttt{REGENERATE.md}: the exact
CaDiCaL invocation for each, with the expected byte count and the SHA-256 the
result must have. Regenerating them
puts a solver back in the loop for the \emph{production} of the proof, which
is where a solver has always been allowed to sit; the check that follows still
does not trust it. The generating pipeline
(\texttt{certify.py}, \texttt{qec\_lib.py}, \texttt{run\_all.sh}) is included
for completeness and is not needed to check anything; it requires
\texttt{numpy} and a CaDiCaL binary, and as noted in D3 its full
\texttt{run\_all.sh} path additionally expects an \texttt{lrat-check} binary
that is not vendored.

\section*{Acknowledgments}

The debt to Bravyi, Cross, Gambetta, Maslov, Rall and Yoder is total: the codes
are theirs, the distances are theirs, and the $\dX=\dZ$ lemma that halves the
work at every BB instance is theirs. To Landahl, Anderson and Rice for the
integer-programming distance method that produced the reference values, and to
Cruz-Benito, Cross, Kremer and Faro for closing the gross-code MILP to gap
zero and thereby fixing the value we certify. To the
LEAN-QEC authors for putting kernel-checked quantum distance proofs on the map,
for naming the gross code as their target in print, and for reaching it in
their repository before we finished writing; several claims in an earlier draft
of this note had to be withdrawn on discovering that, and the discovery was
entirely to the good. To Chen, Jafari and Lai for the SAT benchmark against
which our $n=288$ bound is measured. To Marijn Heule and coauthors for LRAT and
\texttt{drat-trim}, without which none of this would have an interchange
format, and to Armin Biere and coauthors for CaDiCaL. To Carsten Sinz for the
cardinality encoding and to G.~S.~Tseitin for the gate encoding that the whole
construction rests on.

Finding that a result one has produced was produced first by someone else is
the ordinary condition of working in a field that is moving faster than the
writing; we have tried to record it accurately rather than minimally.

\begin{thebibliography}{99}

\bibitem[BCGMRY]{BCGMRY}
S.~Bravyi, A.~W. Cross, J.~M. Gambetta, D.~Maslov, P.~Rall and T.~J. Yoder,
\emph{High-threshold and low-overhead fault-tolerant quantum memory},
Nature \textbf{627} (2024), 778--782; arXiv:2308.07915.

\bibitem[Bie20]{CaDiCaL}
A.~Biere, K.~Fazekas, M.~Fleury and M.~Heisinger,
\emph{CaDiCaL, Kissat, Paracooba, Plingeling and Treengeling entering the SAT
Competition 2020}, in: Proc.\ SAT Competition 2020, Dept.\ of Computer
Science Report Series B, Univ.\ of Helsinki, 51--53.
Version used here: CaDiCaL 3.0.1.

\bibitem[CS96]{CS96}
A.~R. Calderbank and P.~W. Shor,
\emph{Good quantum error-correcting codes exist},
Phys. Rev. A \textbf{54} (1996), 1098--1105.

\bibitem[CJL26]{ChenJafariLai}
Chen, Jafari and Lai,
arXiv:2606.12445, May 29, 2026 (SAT- and MaxSAT-based computation of quantum
code distances). Repository \texttt{guluchen/QDistSAT}; inspected
August 4, 2026, and containing no proof or certificate artifacts.

\bibitem[CCKF26]{CruzBenito}
J.~Cruz-Benito, A.~W. Cross, F.~Kremer and I.~Faro (IBM Quantum),
arXiv:2606.02418, June 1, 2026. Computes the minimum distance of the gross
code $[[144,12,12]]$ exactly by mixed-integer linear programming, closing the
MIP gap to $0$; no optimality certificate is exported.

\bibitem[CFHKS17]{LRAT}
L.~Cruz-Filipe, M.~J.~H. Heule, W.~A. Hunt~Jr., M.~Kaufmann and
P.~Schneider-Kamp,
\emph{Efficient certified RAT verification},
in: Automated Deduction --- CADE-26, Lecture Notes in Comput. Sci.\
\textbf{10395}, Springer, 2017, 220--236.

\bibitem[ELWT26]{LeanQEC}
Ehatamm, Lee, Wu and Tao,
arXiv:2605.16523v1, May 15, 2026 (the LEAN-QEC system paper: SAT-based
quantum code distance proofs replayed in the Lean~4 kernel). Repository
\texttt{VerifiedQC/Lean-QEC}; statements about the repository in this note
refer to commit \texttt{c73827d} of July 10, 2026, inspected August 4, 2026.
Only v1 of the preprint existed at that date.

\bibitem[Gid21]{Stim}
C.~Gidney,
\emph{Stim: a fast stabilizer circuit simulator},
Quantum \textbf{5} (2021), 497. The documentation of
\texttt{search\_for\_undetectable\_logical\_errors} states that the method is
heuristic.

\bibitem[Got97]{Gottesman97}
D.~Gottesman,
\emph{Stabilizer codes and quantum error correction},
Ph.D. thesis, California Institute of Technology, 1997;
arXiv:quant-ph/9705052.

\bibitem[HHW13]{DRATtrim}
M.~J.~H. Heule, W.~A. Hunt~Jr.\ and N.~Wetzler,
\emph{Trimming while checking clausal proofs},
in: Formal Methods in Computer-Aided Design (FMCAD), 2013, 181--188.

\bibitem[LAR11]{LandahlAndersonRice}
A.~J. Landahl, J.~T. Anderson and P.~R. Rice,
\emph{Fault-tolerant quantum computing with color codes},
arXiv:1108.5738, 2011. Source of the mixed-integer-programming distance
computation used in \cite{BCGMRY}.

\bibitem[PSKK22]{QDistRnd}
L.~P. Pryadko, V.~A. Shabashov and V.~K. Kozin,
\emph{QDistRnd: A GAP package for computing the distance of quantum
error-correcting codes},
J. Open Source Softw.\ \textbf{7} (2022), 4120;
\texttt{doi:10.21105/joss.04120}. Upper bounds only, with no performance
guarantee.

\bibitem[Sin05]{Sinz05}
C.~Sinz,
\emph{Towards an optimal CNF encoding of Boolean cardinality constraints},
in: Principles and Practice of Constraint Programming --- CP 2005, Lecture
Notes in Comput. Sci.\ \textbf{3709}, Springer, 2005, 827--831.

\bibitem[Ste96]{Steane96}
A.~M. Steane,
\emph{Error correcting codes in quantum theory},
Phys. Rev. Lett.\ \textbf{77} (1996), 793--797.

\bibitem[Sze26]{PBLean}
S.~Szeider et al.,
arXiv:2602.08692, 2026 (PBLean: importing VeriPB pseudo-Boolean proof
certificates into Lean~4).

\bibitem[Tse68]{Tseitin68}
G.~S. Tseitin,
\emph{On the complexity of derivation in propositional calculus},
in: Studies in Constructive Mathematics and Mathematical Logic, Part II,
1968, 115--125.

\bibitem[Var97]{Vardy97}
A.~Vardy,
\emph{The intractability of computing the minimum distance of a code},
IEEE Trans. Inform. Theory \textbf{43} (1997), 1757--1766.

\bibitem[WJH26]{WJH}
Webster, Jacob and Higgott,
arXiv:2603.22532, March 23, 2026 (a survey and comparison of exact and
heuristic methods for quantum code distance).

\end{thebibliography}

\end{document}
